% from ega4-2-fr.tex
\oldpage[IV-2]{188}

pour un tel idéal $\mathfrak p$, on a nécessairement
$\dim(A/\mathfrak p)=1$; donc a) implique b). Inversement, b) implique que
tout idéal premier $\mathfrak p\in\operatorname{Ass}(A)$ est distinct de
$\mathfrak m$, donc minimal, et par suite b) implique a). On sait déjà que
a) et c) sont équivalentes pour un anneau local de dimension 1 (5.7.8).
Comme il revient au même de dire que $A$ est un anneau de Cohen--Macaulay ou
que $\widehat A$ est un anneau de Cohen--Macaulay (0, 16.5.2), et que l'on
a $\dim(\widehat A)=1$, on voit enfin que $a'$) et $b'$) sont équivalentes à
c).

\begin{proposition}[7.2.2]
\label{IV.2.s7.2.2-fr}
Soient $A$ un anneau local noethérien intègre; posons $A'=\widehat A$,
$X=\operatorname{Spec}(A)$, $X'=\operatorname{Spec}(A')$ et désignons par
$f:X'\to X$ le morphisme canonique. Soient $a$ le point fermé de $X$,
$j:X-\{a\}\to X$ l'injection canonique. Soient $\mathscr F$ un
$\mathscr O_X$-Module cohérent,
$\mathscr F'=f^*(\mathscr F)=\mathscr F\otimes_{\mathscr O_X}\mathscr O_{X'}$;
alors les conditions suivantes sont équivalentes:
\begin{enumerate}
\item[a)] Le $\mathscr O_X$-Module
$j_*(\mathscr F|_{X-\{a\}})$ est cohérent.
\item[b)] Pour tout $x'\in\operatorname{Ass}(\mathscr F')$, on a
$\dim(\overline{\{x'\}})\geqslant2$.
\end{enumerate}
\end{proposition}

\begin{proof}
Soit en effet $a'$ le point fermé de $X'$, qui est l'unique point de la
fibre $f^{-1}(a)$, et soit $j':X'-\{a'\}\to X'$ l'injection canonique.
Comme le morphisme $f$ est fidèlement plat et quasi-compact, il est
équivalent de dire que $j_*(\mathscr F|_{X-\{a\}})$ est cohérent ou que
$j'_*(\mathscr F'|_{X'-\{a'\}})$ est cohérent (5.9.5). D'autre part, comme
$A'$ est isomorphe au quotient d'un anneau local régulier par le théorème de
Cohen (0, 19.8.8), il résulte de (5.11.4) que la condition b) de l'énoncé
équivaut au fait que $j'_*(\mathscr F'|_{X'-\{a'\}})$ est cohérent. D'où la
proposition.

Au lieu d'appliquer (5.11.4) en utilisant le th. de Cohen, ce qui (par
(5.11.2)) fait implicitement appel aux résultats cohomologiques du chap.
III, on peut aussi utiliser le fait que $A'$ est universellement caténaire
(5.6.4) et universellement japonais en vertu de (7.6.5) (la démonstration de
ce dernier résultat n'utilisant pas (7.2.2)).
\end{proof}

\begin{proposition}[7.2.3]
\label{IV.2.s7.2.3-fr}
Soit $A$ un anneau local noethérien intègre. Les notations étant celles de
(7.2.2), les conditions suivantes sont équivalentes:
\begin{enumerate}
\item[a)] $A^{(1)}$ est une $A$-algèbre finie.
\item[b)] Pour toute partie fermée $T$ de $X$ de codimension $\geqslant2$,
si l'on désigne par $i:X-T\to X$ l'injection canonique,
$i_*(\mathscr O_{X-T})$ est un $\mathscr O_X$-Module cohérent.
\item[c)] Pour tout $x'\in\operatorname{Ass}(\mathscr O_{X'})$ et toute
partie fermée $T$ de $X$ de codimension $\geqslant2$, on a
$\operatorname{codim}(f^{-1}(T)\cap\overline{\{x'\}},
\overline{\{x'\}})\geqslant2$.
\end{enumerate}
\end{proposition}

\begin{proof}
Posons $Z=Z^{(2)}(X)$ (5.10.13) et $Z'=f^{-1}(Z)$, qui sont des parties
stables par spécialisation; les conditions a) et b) équivalent
respectivement aux propriétés suivantes:

$a_1)$ $\mathscr H^0_{X/Z}(\mathscr O_X)$ est cohérent;

$b_1)$ $\mathscr H^0_{X/T}(\mathscr O_X)$ est cohérent pour toute partie
fermée $T$ de $X$ de codimension $\geqslant2$.

Compte tenu de (5.9.5), ces deux dernières propriétés équivalent
respectivement à:

$a'_1)$ $\mathscr H^0_{X'/Z'}(\mathscr O_{X'})$ est cohérent;

$b'_1)$ posant $T'=f^{-1}(T)$,
$\mathscr H^0_{X'/T'}(\mathscr O_{X'})$ est cohérent pour toute partie
fermée $T$ de $X$ de codimension $\geqslant2$.

Or, tout point de $\operatorname{Ass}(\mathscr O_{X'})$ se projette dans
$X$ sur le point générique de $X$ puisque $f$ est un morphisme plat
(3.3.2); comme par définition ce point n'appartient pas à $Z$, on voit que
$\operatorname{Ass}(\mathscr O_{X'})$ ne rencontre pas $Z'$;
l'équivalence de $a'_1)$ et $b'_1)$ résulte donc de (5.11.5), puisque $A'$
est isomorphe à un quotient d'anneau régulier en vertu du théorème de Cohen
(0, 19.8.8) (ou encore

